% =============================================================================
%  SECCIONES REDACTADAS — PARTE 3
%  Capítulo 5 (Discusión): secciones 5.2 a 5.6
%  Capítulo 6 (Conclusiones): completo
% =============================================================================


% ╔═══════════════════════════════════════════════════════════════════════════╗
% ║  SUSTITUYE LA SECCIÓN 5.2 COMPLETA, INCLUIDA LA TABLA 5.1                 ║
% ╚═══════════════════════════════════════════════════════════════════════════╝

\section{Comparativa de rendimiento de los modelos}

La Tabla~\ref{tab:resultados} recoge el rendimiento de las diez formulaciones
evaluadas, ordenadas por error mediano. Todas comparten los mismos catorce
orígenes de validación, el mismo horizonte y las mismas métricas.

\begin{table}[H]
\centering
\caption{Comparativa de rendimiento de los modelos evaluados}
\label{tab:resultados}
\small
\begin{tabular}{|l|r|r|r|r|r|r|}
\hline
\rowcolor{grisclaro}
\textbf{Modelo} & \textbf{MAPE} & \textbf{MAPE} & \textbf{Desv.} &
\textbf{Peor} & \textbf{MAE} & \textbf{Frente al} \\
\rowcolor{grisclaro}
 & \textbf{medio} & \textbf{mediana} & \textbf{típica} & \textbf{caso} &
\textbf{(MW)} & \textbf{ingenuo} \\
\hline
\multicolumn{7}{|l|}{\textit{Aprendizaje automático}} \\
\hline
XGBoost          & 1,65 & \textbf{1,34} & 0,70 & 3,10 & 470 & $+$56,2\,\% \\
\hline
Random Forest    & 1,86 & 1,47 & 0,98 & 3,98 & 520 & $+$51,9\,\% \\
\hline
LightGBM         & 1,67 & 1,53 & 0,75 & 3,71 & 473 & $+$50,0\,\% \\
\hline
Ridge            & 2,55 & 1,94 & 1,56 & 6,22 & 717 & $+$36,6\,\% \\
\hline
\multicolumn{7}{|l|}{\textit{Línea base estadística}} \\
\hline
SARIMA clásico   & 3,95 & 2,42 & 2,98 & 10,29 & 1.120 & $+$20,9\,\% \\
\hline
Prophet          & 4,86 & 4,78 & 1,28 & 7,34 & 1.370 & $-$56,4\,\% \\
\hline
SARIMAX + Fourier & 5,66 & 4,90 & 3,18 & 13,02 & 1.565 & $-$60,3\,\% \\
\hline
\multicolumn{7}{|l|}{\textit{Modelos ingenuos}} \\
\hline
Ingenuo semanal  & 3,92 & 3,06 & 2,64 & 10,35 & 1.108 & referencia \\
\hline
Ingenuo diario   & 4,98 & 3,21 & 4,57 & 15,36 & 1.387 & $-$4,9\,\% \\
\hline
Perfil semanal   & 8,54 & 8,30 & 3,33 & 13,60 & 2.434 & $-$171,2\,\% \\
\hline
\end{tabular}
\end{table}

\subsection*{Resultado principal}

El modelo XGBoost alcanza un error mediano del \textbf{1,34\,\%}, lo que supone
una mejora del \textbf{56,2\,\%} sobre la mejor línea base ingenua. En términos
absolutos, el error absoluto medio de 470\,MW representa un 1,7\,\% sobre una
demanda media de 28.263\,MW.

Ambas cifras superan holgadamente los umbrales establecidos en los objetivos del
trabajo —MAPE inferior al 8\,\% y MAE por debajo del 5\,\% del consumo medio— y
se sitúan en el orden de magnitud que maneja el propio operador del sistema en
sus previsiones a día vista.

\subsection*{Los modelos estadísticos no superan la línea base}

Un resultado que merece reportarse con la misma claridad es que \textbf{Prophet y
SARIMAX con términos de Fourier obtienen peores resultados que la predicción
ingenua semanal}. Únicamente el SARIMA clásico logra superarla, con una mejora del
20,9\,\%.

Este hallazgo no constituye una anomalía. La predicción ingenua estacional es una
referencia exigente en demanda eléctrica, precisamente porque la serie presenta
una estructura semanal muy marcada. La observación de
\textcite{hong2016probabilistic} sobre la ausencia habitual de comparación con
líneas base cobra aquí pleno sentido: un trabajo que hubiese omitido estos
modelos ingenuos habría presentado a Prophet como una aportación valiosa cuando
en realidad no mejora la alternativa trivial.

En el caso de SARIMAX, el análisis por origen revela que sus dos peores
resultados corresponden a domingos. La causa es identificable: con dos armónicos
para el periodo de ciento sesenta y ocho horas, la forma semanal representable es
una sinusoide suave, incapaz de reproducir la transición abrupta entre días
laborables y fin de semana.

\subsection*{Estabilidad frente a precisión}

La columna de desviación típica introduce un matiz relevante. El modelo ingenuo
diario presenta una mediana aceptable (3,21\,\%) pero la mayor dispersión de todo
el conjunto (4,57), con un peor caso del 15,36\,\%. El análisis por origen muestra
que sus fallos se concentran en los lunes, con errores del 15,36\,\% y el
13,94\,\%, frente a valores próximos al 2\,\% en el resto. La explicación es
inmediata: predice el lunes a partir del domingo.

Se trata de un modelo que funciona razonablemente bien la mayor parte del tiempo y
falla de forma severa en circunstancias identificables. En un contexto de
explotación, esa característica resulta más problemática que un error medio
ligeramente superior pero estable.

El contraste con Prophet ilustra el compromiso. Prophet obtiene el peor error
mediano de las formulaciones estadísticas (4,78\,\%), pero la menor desviación
típica con diferencia (1,28): nunca destaca y nunca falla estrepitosamente.

\subsection*{Control del sobreajuste}

La Tabla~\ref{tab:sobreajuste} recoge la comparación entre el error de
entrenamiento y el de producción para los modelos de aprendizaje automático, donde
ambos se calculan sobre el mismo horizonte y la comparación resulta por tanto
interpretable.

\begin{table}[H]
\centering
\caption{Error de entrenamiento frente a error de producción}
\label{tab:sobreajuste}
\begin{tabular}{|l|r|r|r|}
\hline
\rowcolor{grisclaro}
\textbf{Modelo} & \textbf{Entrenamiento} & \textbf{Producción} & \textbf{Diferencia} \\
\hline
XGBoost       & 0,67\,\% & 1,65\,\% & $+$146\,\% \\
\hline
LightGBM      & 0,83\,\% & 1,67\,\% & $+$101\,\% \\
\hline
Random Forest & 1,38\,\% & 1,86\,\% & $+$35\,\% \\
\hline
Ridge         & 2,99\,\% & 2,55\,\% & $-$15\,\% \\
\hline
\end{tabular}
\end{table}

XGBoost presenta la mayor diferencia relativa, consecuencia esperable de su
configuración con cuatrocientos árboles de profundidad siete. Random Forest exhibe
la menor, coherente con su parametrización más restrictiva —profundidad doce y
diez observaciones mínimas por hoja—, que penaliza la precisión a cambio de una
mayor capacidad de generalización.

El dato de Ridge requiere interpretación específica: su error de producción es
\textbf{inferior} al de entrenamiento. Ello no responde a ninguna singularidad del
modelo, sino a que \textbf{el periodo de evaluación resulta más benigno que la
media del histórico de entrenamiento}. Las catorce ventanas se sitúan entre
octubre y diciembre de 2024, tramo relativamente estable, mientras que el
entrenamiento abarca tres años que incluyen la volatilidad de 2022.

La implicación afecta a la lectura del resto de la tabla: si el conjunto de
evaluación es más favorable que la media, las diferencias de los modelos de
árboles están \textbf{infravaloradas}. En un periodo de mayor variabilidad cabría
esperar una separación superior. Esta circunstancia se recoge entre las
limitaciones del trabajo.

\subsection*{Selección del modelo}

XGBoost obtiene el mejor resultado en error mediano, en error absoluto medio y en
la evaluación sobre días festivos, por lo que se selecciona como modelo
definitivo.

Conviene no obstante dejar constancia de que LightGBM constituye una alternativa
defendible: su rendimiento en producción es prácticamente equivalente (1,67\,\%
frente a 1,65\,\% de error medio), presenta la mitad de diferencia entre
entrenamiento y producción, y su tiempo de ajuste es sensiblemente menor. En un
escenario que exigiese reentrenamiento frecuente, esas propiedades podrían
inclinar la elección.


% ╔═══════════════════════════════════════════════════════════════════════════╗
% ║  SUSTITUYE LA SECCIÓN 5.3                                                 ║
% ╚═══════════════════════════════════════════════════════════════════════════╝

\section{Análisis de la importancia de variables}

% ─────────────────────────────────────────────────────────────────────────────
% PENDIENTE: incorporar la figura 15_importancia_variables.png y la figura
% 16_shap_summary.png generadas por el notebook 03, junto con el ranking
% concreto de variables que devuelve la sección 6 de dicho cuaderno.
% ─────────────────────────────────────────────────────────────────────────────

\begin{figure}[H]
\centering
\includegraphics[width=0.85\textwidth]{figuras/15_importancia_variables.png}
\caption{Importancia de las variables en el modelo seleccionado}
\label{fig:importancia}
\end{figure}

El análisis de importancia responde al objetivo específico OE6, relativo a la
identificación de los factores con mayor influencia sobre la demanda.

% COMPLETAR con el ranking concreto. Redacción orientativa:
%
% Las variables de mayor peso corresponden a los desfases de 168 y 24 horas,
% seguidas de las variables de calendario. Este orden resulta coherente con los
% hallazgos del análisis exploratorio: el periodograma identificó los ciclos
% semanal y diario como los dominantes, y la clasificación de anomalías situó en
% los festivos cerca del 70 % de las desviaciones significativas.

\subsection*{Valores SHAP}

La importancia nativa de los modelos de árboles indica con qué frecuencia se
emplea cada variable, pero no en qué dirección ni bajo qué condiciones. Los
valores SHAP \parencite{lundberg2017shap} descomponen cada predicción individual
en la contribución aditiva de cada predictor, lo que permite un análisis más
matizado.

\begin{figure}[H]
\centering
\includegraphics[width=0.85\textwidth]{figuras/16_shap_summary.png}
\caption{Valores SHAP del modelo seleccionado}
\label{fig:shap}
\end{figure}

% COMPLETAR con la interpretación de la figura.


% ╔═══════════════════════════════════════════════════════════════════════════╗
% ║  SUSTITUYE LA SECCIÓN 5.4                                                 ║
% ╚═══════════════════════════════════════════════════════════════════════════╝

\section{Análisis de errores}

Las métricas agregadas ocultan la distribución del error. Esta sección examina
dónde falla el modelo, con especial atención a si corrige las deficiencias
detectadas en la línea base.

\subsection*{Error según el tipo de día}

La Tabla~\ref{tab:tipodia} recoge el error desagregado entre días laborables y
fines de semana.

\begin{table}[H]
\centering
\caption{Error según el tipo de día (modelos de aprendizaje automático)}
\label{tab:tipodia}
\begin{tabular}{|l|r|r|r|}
\hline
\rowcolor{grisclaro}
\textbf{Modelo} & \textbf{Laborable} & \textbf{Fin de semana} & \textbf{Diferencia} \\
\hline
XGBoost       & 1,53\,\% & 1,87\,\% & 0,34 \\
\hline
LightGBM      & 1,47\,\% & 1,99\,\% & 0,52 \\
\hline
Random Forest & 1,67\,\% & 2,55\,\% & 0,88 \\
\hline
Ridge         & 2,47\,\% & 2,76\,\% & 0,29 \\
\hline
\end{tabular}
\end{table}

La diferencia entre ambos tipos de día resulta reducida en los modelos de
potenciación del gradiente, lo que indica que las variables de calendario cumplen
su función.

\subsection*{Corrección de las deficiencias de la línea base}

El análisis por origen de la fase anterior había identificado dos patrones de
fallo: el modelo ingenuo diario se degradaba los lunes y SARIMAX los domingos. La
Tabla~\ref{tab:correccion} compara el comportamiento en esos mismos orígenes.

\begin{table}[H]
\centering
\caption{Error en los orígenes problemáticos de la línea base}
\label{tab:correccion}
\begin{tabular}{|l|l|r|r|}
\hline
\rowcolor{grisclaro}
\textbf{Origen} & \textbf{Modelo deficiente} & \textbf{Línea base} & \textbf{XGBoost} \\
\hline
Lunes 28-oct   & Ingenuo diario  & 15,36\,\% & 2,75\,\% \\
\hline
Lunes 2-dic    & Ingenuo diario  & 13,94\,\% & 1,48\,\% \\
\hline
Domingo 17-nov & SARIMAX         & 11,18\,\% & 1,08\,\% \\
\hline
Domingo 22-dic & SARIMAX         & 13,02\,\% & 1,24\,\% \\
\hline
Viernes 27-dic & Ingenuo semanal & 10,35\,\% & 2,01\,\% \\
\hline
\end{tabular}
\end{table}

La corrección es sistemática y de magnitud considerable. Confirma la hipótesis
formulada al analizar la línea base: los fallos identificados obedecían a la
ausencia de información de calendario, y los modelos que la incorporan los
resuelven.

\subsection*{Comportamiento lineal frente a no lineal}

Un resultado con valor explicativo se obtiene al comparar Ridge con los modelos de
árboles en esos mismos orígenes. Ridge mantiene errores del 6,30\,\% y el 5,10\,\%
en los lunes problemáticos, y del 4,06\,\% en uno de los domingos, mientras las
tres formulaciones basadas en árboles se sitúan por debajo del 3\,\%.

La circunstancia resulta ilustrativa porque \textbf{Ridge dispone exactamente de
las mismas variables}, incluida la indicadora de día laborable. Su limitación es
estructural: al ser lineal, únicamente puede sumar un efecto constante asociado al
tipo de día, sin capturar que \textbf{dicho efecto depende de la hora}. Un lunes a
las cuatro de la madrugada se asemeja considerablemente a un domingo, mientras que
a las nueve de la mañana la diferencia es máxima. Esa interacción entre hora y día
de la semana solo la reproducen los modelos capaces de representar relaciones no
lineales.

\subsection*{Distribución del error a lo largo del día}

El error presenta un mínimo en torno a las cinco de la madrugada (0,96\,\%) y un
máximo a las diecisiete horas (2,63\,\%). El patrón es coherente con la estructura
del problema: las horas de rampa —el ascenso matinal y el pico vespertino—
concentran la mayor incertidumbre, mientras que el valle nocturno resulta
altamente predecible.

\subsection*{Calidad global del ajuste}

El coeficiente de correlación entre valores observados y predichos alcanza 0,9931,
con un error máximo de 1.948\,MW sobre una demanda media de 28.263\,MW. El sesgo
medio es de $-204$\,MW, lo que indica una ligera tendencia a la infraestimación sin
relevancia práctica.


% ╔═══════════════════════════════════════════════════════════════════════════╗
% ║  NUEVA SECCIÓN 5.5 — INSERTAR ANTES DE "Aplicación desarrollada"          ║
% ╚═══════════════════════════════════════════════════════════════════════════╝

\section{Evaluación en días festivos}
\label{sec:festivos}

El análisis exploratorio estableció que cerca del 70\,\% de las anomalías se
concentraba en días festivos, y sobre esa base se incorporaron las variables de
calendario. Procede verificar de forma explícita que dichas variables cumplen su
cometido, para lo cual se emplea el conjunto de evaluación dirigido descrito en la
sección~\ref{sec:validacion}.

\begin{table}[H]
\centering
\caption{Rendimiento en los ocho días festivos evaluados}
\label{tab:festivos}
\begin{tabular}{|l|r|r|r|r|}
\hline
\rowcolor{grisclaro}
\textbf{Modelo} & \textbf{MAPE en} & \textbf{Peor} & \textbf{MAPE} &
\textbf{Penali-} \\
\rowcolor{grisclaro}
 & \textbf{festivos} & \textbf{festivo} & \textbf{ordinario} & \textbf{zación} \\
\hline
XGBoost         & 3,32\,\% & 7,02\,\% & 1,65\,\% & 2,01 \\
\hline
Random Forest   & 4,10\,\% & 8,31\,\% & 1,86\,\% & 2,21 \\
\hline
Ingenuo diario  & 14,12\,\% & 19,52\,\% & 4,98\,\% & 2,84 \\
\hline
Ingenuo semanal & 16,63\,\% & 40,68\,\% & 3,92\,\% & 4,24 \\
\hline
\end{tabular}
\end{table}

La columna de penalización expresa el cociente entre el error en festivo y el
error en día ordinario. Un valor próximo a la unidad indicaría que el modelo
aprovecha íntegramente la información de calendario.

El caso más ilustrativo corresponde al 25 de diciembre: el modelo ingenuo semanal
alcanza un error del \textbf{40,68\,\%}, frente al \textbf{4,74\,\%} de XGBoost.
La diferencia es de casi un orden de magnitud y responde a una causa evidente: el
modelo ingenuo compara la Navidad con el miércoles precedente, mientras que el
modelo de aprendizaje automático dispone de la indicación de festividad.

\subsection*{Festivos nacionales y autonómicos}

Una primera versión de este análisis reveló un patrón inesperado. Los modelos de
aprendizaje automático superaban ampliamente a la línea base en los festivos
nacionales, pero no en los autonómicos, donde el error resultaba comparable e
incluso superior.

La causa admite una explicación precisa. La serie de demanda es
\textbf{peninsular}, mientras que el calendario empleado correspondía a la
\textbf{Comunidad de Madrid}. El 2 de mayo, festividad exclusivamente madrileña,
el modelo recibía la indicación de festivo y anticipaba una caída pronunciada del
consumo; sin embargo, el resto del país mantenía su actividad ordinaria y la
demanda apenas descendía.

La solución adoptada sustituye la variable binaria por una \textbf{intensidad
ponderada por población}: la fracción de habitantes del territorio nacional para
los que cada fecha constituye día festivo. Se consulta a tal efecto el calendario
de las diecinueve subdivisiones y se pondera por población.

\begin{table}[H]
\centering
\caption{Intensidad de festivo según el alcance territorial}
\label{tab:intensidad}
\begin{tabular}{|l|l|c|c|}
\hline
\rowcolor{grisclaro}
\textbf{Fecha} & \textbf{Festividad} & \textbf{Comunidades} & \textbf{Intensidad} \\
\hline
25 de diciembre & Natividad del Señor & 19 & 1,000 \\
\hline
28 de marzo & Jueves Santo & 17 & 0,725 \\
\hline
25 de julio & Santiago Apóstol & 5 & 0,271 \\
\hline
2 de mayo & Fiesta de la Comunidad de Madrid & 1 & 0,143 \\
\hline
\end{tabular}
\end{table}

Resulta pertinente señalar que una simple distinción binaria entre festivos
nacionales y autonómicos tampoco habría resuelto el problema. El Jueves Santo
figura como autonómico en las fuentes de calendario, pero afecta en realidad al
72,5\,\% de la población. La representación mediante una magnitud continua es la
única que recoge adecuadamente este matiz.

\subsection*{Efecto de la corrección}

El resultado sobre los dos casos problemáticos fue sustancial. El error del 2 de
mayo descendió del \textbf{10,62\,\% al 2,14\,\%}, y el de Santiago Apóstol del
4,28\,\% al 2,89\,\%. En términos agregados, el error medio en festivos pasó del
4,81\,\% al 3,32\,\%, con una reducción de la penalización de 2,95 a 2,01.

El rendimiento global permaneció prácticamente inalterado, como cabía esperar: las
festividades representan alrededor de veinte días al año y su peso sobre la
métrica media es reducido. Una corrección que mejora los casos difíciles sin
degradar el caso general constituye el resultado deseable.


% ╔═══════════════════════════════════════════════════════════════════════════╗
% ║  SUSTITUYE LA SECCIÓN 5.5 ACTUAL, QUE PASA A SER 5.6                      ║
% ╚═══════════════════════════════════════════════════════════════════════════╝

\section{Aplicación desarrollada}

% ─────────────────────────────────────────────────────────────────────────────
% PENDIENTE: incorporar una captura de la aplicación en funcionamiento.
% ─────────────────────────────────────────────────────────────────────────────

\begin{figure}[H]
\centering
% \includegraphics[width=\textwidth]{figuras/22_captura_aplicacion.png}
\framebox[0.9\textwidth][c]{\rule{0pt}{6cm}\textit{[Captura de la aplicación]}}
\caption{Interfaz de la aplicación web desarrollada}
\label{fig:app}
\end{figure}

La aplicación materializa el objetivo específico OE7. Su elemento central es la
traducción de la predicción numérica a una recomendación en lenguaje natural,
generada automáticamente a partir de las franjas horarias de menor demanda
prevista.

\subsection*{Degradación del esquema recursivo}

El modelo opera con un horizonte nativo de veinticuatro horas. La extensión a
cuarenta y ocho mediante realimentación de las predicciones introduce una
acumulación de error que conviene cuantificar.

% COMPLETAR con los valores que devuelve la celda 6 del notebook 04:
%   MAPE horas  1–24: __ %
%   MAPE horas 25–48: __ %

La aplicación advierte de esta circunstancia cuando el usuario selecciona el
horizonte extendido.

\subsection*{Comportamiento ante fallos de las fuentes}

Durante el desarrollo se constató que la interfaz REData devuelve errores de
servidor ante peticiones que abarcan periodos extensos, dado que publica con
resolución de diez minutos. La descarga se fragmentó en consecuencia en bloques de
cinco días, con tolerancia a fallos parciales: si un bloque concreto no responde,
los restantes se procesan y la aplicación advierte de la posible existencia de
huecos.

Esta consideración excede el ámbito estrictamente académico, pero resulta
determinante para un sistema destinado a funcionar de forma continuada.


% ╔═══════════════════════════════════════════════════════════════════════════╗
% ║  SUSTITUYE EL CAPÍTULO 6 COMPLETO                                         ║
% ╚═══════════════════════════════════════════════════════════════════════════╝

\chapter{Conclusiones}

\section{Cumplimiento de los objetivos}

Se revisa a continuación el grado de consecución de cada objetivo específico.

\paragraph{OE1 — Análisis del comportamiento histórico.}
\textbf{Alcanzado.} El análisis exploratorio caracterizó la serie e identificó las
tres escalas de periodicidad presentes. El tratamiento de valores ausentes reveló
además una circunstancia metodológicamente relevante: las ausencias no figuraban
como valores nulos sino como registros inexistentes en el índice temporal, de modo
que los procedimientos habituales de interpolación no habrían producido efecto
alguno.

\paragraph{OE2 — Incorporación de variables exógenas.}
\textbf{Alcanzado.} Se integraron las variables meteorológicas de AEMET y el
calendario laboral. Este último evolucionó durante el trabajo desde una variable
binaria hasta una intensidad ponderada por población, tras detectarse la
discordancia entre el ámbito territorial de la demanda y el del calendario.

\paragraph{OE3 — Aplicación de series de Fourier.}
\textbf{Alcanzado.} El análisis espectral identificó los periodos de 24, 12 y 168
horas, resultado que fundamenta empíricamente la selección de los desfases
temporales empleados como predictores.

\paragraph{OE4 — Desarrollo y comparación de modelos.}
\textbf{Alcanzado y superado.} Se evaluaron diez formulaciones frente a las seis
inicialmente previstas, incorporando tres modelos ingenuos como línea base
adicional.

\paragraph{OE5 — Evaluación con control del sobreajuste.}
\textbf{Alcanzado.} Se registraron las métricas dentro y fuera de muestra para
todos los modelos, con la precisión relativa a su distinta interpretabilidad en
formulaciones estadísticas y de aprendizaje automático.

\paragraph{OE6 — Identificación de los factores determinantes.}
\textbf{Alcanzado.} Mediante importancia nativa y valores SHAP. El resultado
converge con el del análisis exploratorio, lo que refuerza ambas vías.

\paragraph{OE7 — Despliegue de la aplicación.}
\textbf{Alcanzado.} La aplicación se encuentra desplegada y accesible
públicamente, con predicción a 24 y 48 horas, identificación de franjas favorables
y recomendaciones en lenguaje natural.

\subsection*{Cumplimiento de las métricas objetivo}

\begin{table}[H]
\centering
\caption{Grado de consecución de las métricas establecidas}
\label{tab:objetivos}
\begin{tabular}{|l|c|c|c|}
\hline
\rowcolor{grisclaro}
\textbf{Métrica} & \textbf{Objetivo} & \textbf{Obtenido} & \textbf{Estado} \\
\hline
MAPE & $<$ 8\,\% & 1,34\,\% & Superado \\
\hline
MAE sobre consumo medio & $<$ 5\,\% & 1,7\,\% & Superado \\
\hline
Mejora sobre línea base & Superarla & $+$56,2\,\% & Superado \\
\hline
Aplicación desplegada & 1 URL pública & Sí & Cumplido \\
\hline
\end{tabular}
\end{table}


\section{Principales aportaciones}

\subsection*{Un sistema completo, no un experimento}

La aportación principal no reside en el valor concreto del error alcanzado, sino
en haber recorrido el ciclo íntegro desde el dato abierto hasta la aplicación
accesible. Ese recorrido obligó a resolver cuestiones que un estudio puramente
comparativo no llega a plantearse: la disponibilidad efectiva de cada variable en
el instante de predicción, la coherencia entre el código de entrenamiento y el de
explotación, o el comportamiento del sistema ante el fallo de una fuente externa.

\subsection*{Tratamiento riguroso y documentado de los datos}

Tres decisiones metodológicas merecen destacarse por su carácter no evidente:

\begin{enumerate}
  \item \textbf{El reindexado previo a la imputación.} Las ausencias se
        manifestaban como registros inexistentes, no como valores nulos. Sin
        reindexar, la interpolación no habría producido efecto y la verificación
        posterior habría comparado la serie consigo misma.

  \item \textbf{La estrategia escalonada de imputación.} El desplazamiento
        semanal para huecos extensos preserva el perfil diario, que una
        interpolación lineal destruiría.

  \item \textbf{La detección de anomalías condicionada al contexto.} La regla del
        rango intercuartílico no identificaba ningún valor atípico; el
        procedimiento sobre el residuo respecto al perfil temporal detectó 555.
\end{enumerate}

\subsection*{Cuantificación de la fuga de información}

El trabajo no se limita a evitar la fuga de información, sino que mide su efecto:
la inclusión de variables no disponibles en el instante de predicción reduce el
error aparente en un 59\,\%. Esta cuantificación documenta una decisión
metodológica que la mayoría de los trabajos publicados no llega a explicitar.

\subsection*{Validación de la hipótesis sobre el efecto del calendario}

La cadena argumental del trabajo se cierra de forma completa. El análisis
exploratorio identificó que cerca del 70\,\% de las anomalías se concentraba en
festivos; ese hallazgo motivó la incorporación de las variables de calendario; y
la evaluación dirigida verificó que dichas variables cumplen su cometido, con
diferencias de hasta un orden de magnitud frente a los modelos que carecen de
ellas.

\subsection*{Detección y corrección de una limitación propia}

La discordancia entre el ámbito territorial de la demanda —peninsular— y el del
calendario —autonómico— se detectó al analizar los resultados, no se conocía de
antemano. La solución adoptada, una intensidad ponderada por población, corrigió
el error del 2 de mayo del 10,62\,\% al 2,14\,\% sin degradar el rendimiento
general.


\section{Limitaciones del trabajo}

\paragraph{Resolución de las variables meteorológicas.}
Los valores climatológicos de AEMET tienen periodicidad diaria, mientras que la
variable objetivo es horaria. La propagación del valor diario a las veinticuatro
horas del día renuncia al ciclo térmico intradiario.

\paragraph{Representatividad de la estación meteorológica.}
Se emplea una única estación —Madrid-Retiro— como aproximación de las condiciones
peninsulares. La heterogeneidad climática del territorio nacional no queda
recogida.

\paragraph{Periodo de evaluación acotado.}
Los catorce orígenes cubren aproximadamente setenta días, entre octubre y
diciembre de 2024. El signo negativo de la diferencia entre entrenamiento y
producción en el modelo Ridge sugiere que dicho periodo resulta más estable que la
media del histórico, por lo que las cifras de sobreajuste podrían estar
infravaloradas.

\paragraph{Ausencia de información sobre generación.}
El sistema no incorpora datos de generación renovable, factor que influye
sensiblemente sobre el precio y que podría mejorar la calidad de las
recomendaciones económicas.

\paragraph{Extensión recursiva del horizonte.}
La predicción más allá de veinticuatro horas se obtiene realimentando las
predicciones anteriores, procedimiento que acumula error. Un modelo entrenado
específicamente para cuarenta y ocho horas obtendría mejores resultados.

\paragraph{Relación entre demanda y precio.}
La aplicación emplea la demanda prevista como indicador aproximado del precio. La
relación entre ambas magnitudes es estrecha pero no determinista: la composición
del parque generador disponible introduce una variabilidad no contemplada.


\section{Líneas de trabajo futuro}

\paragraph{Predicción directa del precio.}
Incorporar el precio del mercado diario como variable objetivo adicional
permitiría ofrecer recomendaciones económicas directas en lugar de aproximadas.

\paragraph{Meteorología de resolución horaria.}
El reanálisis ERA5 del programa Copernicus proporciona temperatura horaria
histórica, lo que eliminaría la limitación principal de las variables
meteorológicas actuales.

\paragraph{Modelos fundacionales para series temporales.}
Han surgido recientemente modelos preentrenados de propósito general para
predicción de series temporales, capaces de operar sin entrenamiento específico
sobre el conjunto de destino. Su evaluación comparativa frente al enfoque adoptado
constituye una línea de continuación natural.

\paragraph{Combinación de modelos.}
La evidencia de las competiciones M \parencite{makridakis2020m4} indica que la
combinación de predicciones mejora sistemáticamente la robustez. Un promedio
ponderado de XGBoost y LightGBM resultaría inmediato de implementar.

\paragraph{Desagregación territorial.}
La disponibilidad de demanda por comunidad autónoma permitiría emplear el
calendario correspondiente a cada territorio, resolviendo de raíz la limitación
que motivó la variable de intensidad ponderada.

\paragraph{Integración del consumo individual.}
La incorporación de los datos de consumo del propio usuario, accesibles a través
de las distribuidoras, permitiría pasar de una recomendación general a una
estimación personalizada del ahorro.
